\subsection{Frame Averaging}\label{sec:bg-frame-averaging}
\begin{definition}
    Given a $G$-space $X$, a \emph{frame} is a set-valued function $\FF : X \to 2^G \sminus \set{\emp}$.
\end{definition}

\begin{definition}
    A frame $\FF$ is \emph{$G$-equivariant} if for all $g \in G$ and $x \in X$,
    \[
        \FF(gx) = g\FF(x) = \set{gh : h \in \FF(x)}
    \]
    or \emph{$G$-invariant} if for all $g \in G$ and $x \in X$,
    \[
        \FF(gx) = \FF(x).
    \]
    Note that this corresponds to our usual definition of equivariance/invariance with $G$ acting on $2^G$ element-wise.
\end{definition}

\begin{definition}
    Let $X,Y$ be $G$-spaces where $G$ acts linearly.
    A frame $\FF$ on is \emph{bounded} over $K \sub X$ if there exists $c > 0$ such that for all $x \in K$,
    \[
        \sup_{g \in \FF(x)} \norm{g \cdot -}_{op}
        = \sup_{g \in \FF(x)} \sup_{y \in Y} \norm{gy}
        \leq c.
    \]
\end{definition}

\begin{definition}
    A frame $\FF$ is \emph{finite} if for all $x \in X$, $\card(\FF(x))$ is finite.
    More generally, the \emph{cardinality} of $\FF$ is $\card(\FF) := \sup_{x \in X} \card(\FF(x))$.
\end{definition}

\begin{proposition}\label{prop:invariant-frame-average}
    (Theorem 1 of \cite{puny2022frameaveraginginvariantequivariant})
    Let $\FF$ be a finite equivariant frame and let $f : X \to Y$ where $Y$ is a vector space.
    Then the function $\II_\FF(f) : X \to Y$ given by
    \[
        \II_\FF (f) (x) := \f{1}{\card(\FF(x))} \sum_{g \in \FF(x)} f(\inv{g} x)
    \]
    is $G$-invariant.
    \begin{proof}
        Let $x \in X$ and $h \in G$.
        Then
        \begin{align*}
            \II_\FF (f) (hx)
            &= \f{1}{\card(\FF(hx))} \sum_{g \in \FF(hx)} f(\inv{g} hx) \\
            &= \f{1}{\card(h\FF(x))} \sum_{g \in h\FF(x)} f(\inv{g} hx) \\
            &= \f{1}{\card(\FF(x))} \sum_{g \in \FF(x)} f(\inv{(hg)} hx) \\
            &= \f{1}{\card(\FF(x))} \sum_{g \in \FF(x)} f(\inv{g} x) \\
            &= \II_\FF (f) (x).
        \end{align*}
    \end{proof}
\end{proposition}

\begin{proposition}\label{prop:equivariant-frame-average}
    (Theorem 2 of \cite{puny2022frameaveraginginvariantequivariant})
    Let $\FF$ be a finite equivariant frame and let $f : X \to Y$ where $Y$ is a vector space.
    Then the function $\EE_\FF (f) : X \to Y$ given by
    \[
        \EE_\FF (f) (x) = \f{1}{\card(\FF(x))} g\sum_{g \in \FF(x)} f(\inv{g} x)
    \]
    is $G$-equivariant.
    \begin{proof}
        Let $x \in X$ and $h \in G$.
        Then
        \begin{align*}
            \EE_\FF (f) (hx)
            &= \f{1}{\card(\FF(hx))} \sum_{g \in \FF(hx)} gf(\inv{g} hx) \\
            &= \f{1}{\card(h\FF(x))} \sum_{g \in h\FF(x)} gf(\inv{g} hx) \\
            &= \f{1}{\card(\FF(x))} \sum_{g \in \FF(x)} gf(\inv{(hg)} hx) \\
            &= \f{1}{\card(\FF(x))} \sum_{g \in \FF(x)} gf(\inv{g} x) \\
            &= \EE_\FF (f) (x).
        \end{align*}
    \end{proof}
\end{proposition}

\begin{proposition}\label{prop:frame-is-orbit-union}
    (Theorem 3 of \cite{puny2022frameaveraginginvariantequivariant})
    Let $\FF$ be an equivariant frame and let $x \in X$.
    Then $\FF(x)$ is a (disjoint) union of orbits $G_x g \in \FF(x) / G_x$.
    \begin{proof}
        Observe that it suffices to show that $hg \in \FF(x)$ for all $g \in \FF(x)$ and $h \in G_x$.
        So, fix some $h \in G_x$ and we conclude that $hg \in h\FF(x) = \FF(hx) = \FF(x)$.
    \end{proof}
\end{proposition}

\begin{corollary}
    Given an equivariant frame $\FF$ and function $f : X \to Y$, we can instead compute the averages in propositions \ref{prop:invariant-frame-average} and \ref{prop:equivariant-frame-average} as
    \[
        \II_\FF (f) (x) = \f{\card(G_x)}{\card(\FF(x))} \sum_{G_x g \in \FF(x) / G_x} f(\inv{g}x)
    \]
    and
    \[
        \EE_\FF (f) (x) = \f{\card(G_x)}{\card(\FF(x))} \sum_{G_x g \in \FF(x) / G_x} gf(\inv{g}x)
    \]
    respectively.
\end{corollary}

\begin{corollary}\label{cor:frame-card-lower-bound}
    Given an equivariant frame $\FF$ and $x \in X$, $\card(\FF(x)) \geq \card(G_x)$.
    In particular, the cardinality of $\FF$ is at least $\sup_{x \in X} \card(G_x)$.
\end{corollary}

\begin{remark}\label{rem:frame-card-upper-bound}
    Clearly the largest possible equivariant frame is $\FF(x) = G$ for all $x \in X$, so the cardinality of $\FF$ is at most $\card{G}$.
\end{remark}


\begin{theorem}
    (Theorem 2.8 of \cite{dym2024equivariant})
    Let $G$ be a finite group acting continuously on a locally connected Hausdorff space $X$, and let $\FF$ be an equivariant frame such that $\II_\FF (f)$ is continuous for every map $f : X \to Y$ where $Y$ is some non-trivial (real) vector space.
    If $X_{free}$ is connected, then $\FF(x) = G$ for all $x \in \clo{X_{free}}$.
    \begin{proof}
        Define $U = \set{x \in X_{free} : e \in \FF(x)}$ and $V = X_{free} \sminus U$.
        Since $G$ is finite and acts freely on $X_{free}$, for all $x \in X_{free}$ we can find connected open sets $U_x \ni x$ such that $gU_x \cap U_x = \emp$ for $g \neq e$.
        Fixing some non-zero vector $y \in Y$, there are then continuous maps $f_x : X_{free} \to Y$ such that $f_x[U_x] = y$ and $f_x[gU_x] = 0$ for $g \neq e$.
        Fixing $x \in X_{free}$, we know that $\II_\FF (f_x)$ is continuous, and additionally for $v \in U_X$ we have
        \begin{align*}
            \II_\FF (f_x) (v)
            &= \f{1}{\card(\FF{v})} \sum_{g \in \FF(v)} f_x(\inv{g} v) \\
            &= \f{1}{\card(\FF{v})} \sum_{g \in \FF(v)} \delta_{e,g} y \\
            &= \f{y}{\card(\FF{v})} \ind_{\FF(v)}(e)
        \end{align*}
        Then if $X \in U$ (i.e. $e \in \FF(x)$) we have that $\II_\FF(f_x)(x) = \f{y}{\card(\FF(x))}$, and consequently $\II_\FF (f_x) [U_x] = \f{y}{\card(\FF(x))}$ since $U_x$ is connected and $\II_\FF (f_x)$ is continuous.
        In particular, this means $e \in \FF{v}$ for all $v \in U_x$, so $U_x \sub U$.
        Then $U = \cup_{x \in U} U_x$, so we have that $U$ is open.
        Similarly, if $x \in V$ then $\II_\FF (f_x) [U_x] = 0$ leading to the fact that $V$ is open as well.

        Since $X_{free}$ is connected, we must have that either $U$ or $V$ is empty.
        To see that $U$ is non-empty, let $x \in X_{free}$ be arbitrary and fix some $g \in \FF(x)$.
        Since $\FF$ is equivariant we then have that $e \in \FF(\inv{g} x)$, so $\inv{g} x \in U$.
        Hence, we conclude that $U = X_{free}$ and $e \in \FF(x)$ for all $x \in X_{free}$.
        
        With this, we can easily show that $\FF(x) = G$ for all $x \in X_{free}$.
        Let $x \in X_{free}$ and $g \in G$, and we know that $e \in \FF(\inv{g} x)$ so $g \in \FF(x)$ by equivariance of $\FF$.
        Hence, $\FF(x) = G$ for all $x \in X_{free}$.
        
        To conclude, let $x \in \clo{X_{free}}$, and we can find a connected open set $U \ni x$ such that $gU \cap U \neq \emp$ if and only if $g \in G_x$.
        Then there is a continuous map $f : X \to Y$ such that $f[U] = y$ and $f[gU] = 0$ if $g \not\in G_x$.
        Then for $v \in U$ we have as before that
        \begin{align*}
            \II_\FF (f) (v)
            &= \f{1}{\card(\FF(v))} \sum_{g \in \FF(v)} f(\inv{g} v) \\
            &= \f{1}{\card(\FF(v))} \sum_{g \in \FF(v)} \ind_{G_x}(g) y \\
            &= \f{\card(G_x \cap \FF(v))}{\card(\FF(v))} y.
        \end{align*}
        Then by continuity and connectivity we have that $\II_\FF (f)$ is constant on $U$.
        Since $x \in \clo{X_{free}}$, let $\net{x} \sub X_{free}$ be a net such that $x_d \to x$.
        There is then some $d' \in D$ such that for all $d \geq d'$, $x_d \in U$, and $e \in \FF(U)$ so $\II_\FF (f) (x_d) = \f{\card(G_x \cap \FF(x_d))}{\card(\FF(x_d))}y = \f{\card(G_x)}{\card(G)}y \neq 0$.
        Then by continuity of $\II_\FF (f)$ we have that $\II_\FF (f) (x) = \lim_{d \in D} \II_\FF (f) (x_d) \neq 0$ so $G_x \cap \FF(x) \neq 0$.
        Letting $g \in G_x \cap \FF(x)$, we have that $e \in \inv{g}\FF(x) = \FF(\inv{g}x) = \FF(x)$.
        This shows that $e \in \FF(x)$ for all $x \in \clo{X_{free}}$.
        Finally, for $x \in \clo{X_{free}}$ and $g \in G$, and we know that $e \in \FF(\inv{g} x)$ so $g \in \FF(x)$ by equivariance of $\FF$.
        Therefore $\FF(x) = G$ for all $x \in \clo{X_{free}}$.
    \end{proof}
\end{theorem}

\begin{corollary}
    (Corrollary 2.9 in \cite{dym2024equivariant})
    Let $d,n > 1$ and consider $\Sigma_n$ acting on $\R^{d \times n}$.
    If $\FF$ is a continuity-preserving equivariant frame, then $\FF(A) = \Sigma_n$ for all $A \in \R^{d \times n}$.
    Note that this provides an alternate proof of theorem \ref{thm:no-Sigman-canon-for-dxn}, since a canonicalization is a frame with cardinality 1.
\end{corollary}

\begin{lemma}
    Let $X$ be a $G$-space such that the quotient $q : X \to X/G$ is a principal $G$-bundle (for instance, $G$ is a compact Lie group acting freely on $X$).
    Then there exists a section $s : X/G \to X$ of $q$ if and only if there exists a $G$-equivariant map $\varphi : X \to_G G$.
    \begin{proof}
        First, assume we have a section $s$ of $q$.
        For two elements $x,y \in X$ in the same $G$ orbit, let $g \in G$ be the unique element such that $gy = x$ (this is well-defined since $G$ acts freely), we then define $x/y = g$.
        Now define the function $\varphi : X \to G$ via $\varphi(x) = x/(s \circ q)(x)$.
        Observe that $\varphi$ is $G$-equivariant since $x \mapsto x$ is clearly equivariant and $(s \circ q)(x)$ is invariant.
        To see that $\varphi$ is continuous, let $U \sub G$ be open.
        Then
        \begin{align*}
            \inv{\varphi}[U]
            &= \set{x \in X : x/(s \circ q)(x) \in U} \\
            &= \set{x \in X : \exists g \in U, g (s \circ q)(x) = x} \\
            &= \set{x \in X : (s \circ q)(x) \in \inv{U} x} \\
            &= \inv{(s \circ q)}[\inv{U}x].
        \end{align*}
        Observe that $\inv{U}$ is open since $\inv{(-)}$ is a homeomorphism on any topological group, then $\inv{U}x$ is open as well since $q$ is a principal $G$-bundle.
        Hence, $\inv{\varphi}[U]$ is open since $s \circ q$ is continuous, so we conclude that $\varphi$ is a continuous equivariant map.
        
        Conversely, assume $\varphi$ is an equivariant map.
        Then we can define $s : X/G \to X$ via $s(Gx) = \inv{\varphi(x)} x$.
        Observe that $s$ is indeed a well-defined map since
        \[
            s(Ggx)
            = \inv{\varphi(gx)}gx
            = \inv{g\varphi(x)}gx
            = \inv{\varphi(x)}\inv{g}gx
            = \inv{\varphi(x)}x
            = s(Gx).
        \]
        Then $s$ is continuous by the continuity of $\varphi$ and $\pinv{-}$ and the definition of the quotient topology, so we are done.
    \end{proof}
\end{lemma}

\begin{lemma}\label{lem:cont-preserving-frame-is-cont}
    Let $X$ be a $G$-space with compact $G$ such that $q : X \to X/G$ is a principal $G$-bundle and let $F : X \to G^n/\Sigma_n$.
    If $\II_F (f) (x) = \f{1}{n} \Sigma_{g \in F(x)} f(\inv{g} x)$ is continuous for all $f : X \to \R$, then $F$ is continuous.
    \begin{proof}
        \Afsoc there is some $x_0 \in X$ such that $F$ is not continuous, so we can find a net $\net{x} \sub X$ with $x_d \to x_0$ such that $F(x_d) \not\to F(x_0)$.
        Since $G$ is compact, without loss of generality we can assume that $F(x_d) \to L \neq F(x_0)$ by passing to a subnet.
        Now since $q$ is a principal $G$-bundle, there exists an open neighborhood $U$ of $q(x_0)$ such that $\inv{q}[U] \cong G \times U$.
        This homeomorphism is demonstrated by a section $s : X/G \to X$ of $q$ so we have the map $x \mapsto h \cdot s(q(x))$ for some unique $h \in G$ (this is continuous and well-defined since $q$ is a principal $G$-bundle).
        We can then find a continuous function $\varphi : G \to \R$ such that $\sum_{g \in L} \varphi(\inv{g}h) \neq \sum_{g \in F(x_0)} \varphi(\inv{g}h)$ where $x_0 = h \cdot s(q(x_0))$.
        Then by Urysohn's lemma there is some $\psi : X \to \R$ such that $\psi(x_0) = 1$ and $\psi[X \sminus \inv{q}[U]] = 0$, so the function $f : X \to \R$ defined by
        \[
            f(x)
            = \begin{cases}
                \psi(x) \varphi(g) &\tif x = g \cdot s(q(x)) \in \inv{q}[U] \\
                0 &\totherwise
            \end{cases}
        \]
        is well-defined and continuous.
        Then by assumption we have that
        \[
            \II_F (f)(x_d)
            \to \II_F (f) (x_0)
            = \f{1}{n} \sum_{g \in F(x_0)} f(\inv{g} x_0)
            = \f{1}{n} \sum_{g \in F(x_0)} \varphi(\inv{g}h).
        \]
        However, the map from $G^n/\Sigma_n$ to $\R$ given by
        \[
            \set{g_1,...,g_n} \mapsto \f{1}{n} \sum_i f(\inv{g_i} x)
        \]
        for any fixed $f : X \to R$ and $x \in X$ is also continuous by the definition of the quotient topology, so we have
        \[
            \II_F (f) (x_d)
            = \f{1}{n} \sum_{g \in F(x_d)} f(\inv{g} x_d)
            \to \f{1}{n} \sum_{g \in L} f(\inv{g} x_0) %TODO: check this
            = \f{1}{n} \sum_{g \in L} \varphi(\inv{g}h).
        \]
        This is a contradiction since these limits are not equal, so we conclude that $F$ must be continuous.
    \end{proof}
\end{lemma}

\begin{theorem}
    (Theorem 2.10 of \cite{dym2024equivariant})
    Let $n \geq 2$ and let $\FF$ be a continuity-preserving equivariant frame for $\SO(2)$ acting on $X := \R^{2 \times n} \sminus 0$.
    Then $\card(\FF) = \infty$.
    \begin{proof}
        \Afsoc that we have some finite continuity-preserving frame $\FF$.
        Without loss of generality, we may assume that $\card(\FF(x)) = \card(\FF) =: N$ for all $x \in X$, since if not we may simply repeat elements.
        This means $\FF$ can be thought of as a function from $X \to \SO(2)^N / \Sigma_N$.
    \end{proof}
\end{theorem}

\begin{definition}
    Let $\mu$ be a Borel measure on a topological group $G$, and let $g \in G$.
    Then the \emph{pushforward measure} $g_* \mu$ is defined for Borel sets $E \sub G$ as $g_* \mu(E) := \mu(\inv{g}E)$.
\end{definition}

\begin{definition}
    Let $\mu$ be a Borel measure on a topological group $G$ which acts on $X$, and let $x \in X$.
    Then the measure $\angles{\mu}_x$ is defined for Borel sets $E \sub G$ by
    \[
        \angles{\mu}_x(E)
        := \int_{g \in G_x} g_* \mu(E) dg
        = \int_{g \in G_x} \mu(\inv{g}E) dg
    \]
    where the integral is with respect to the Haar measure of $G_x$.
\end{definition}

\begin{definition}
    Given a $G$-space $X$, a \emph{weighted frame} is a function $\mu_\bullet : X \to \PP(G)$ where $\PP(G)$ is the space of Borel probability measures on $G$
    For $x \in X$ we write $\mu_x = \mu_\bullet(x)$ and $\Bar{\mu_x} := \angles{\mu_x}_x$.
\end{definition}

\begin{definition}
    A weighted frame $\mu_\bullet$ is \emph{equivariant} at $x \in X$ if for all $g \in G$, $\mu_{gx} = g_* \mu_x$.
    $\mu_\bullet$ is \emph{equivariant} if it is equivariant at all $x \in X$.
\end{definition}

\begin{definition}
    A weighted frame $\mu_\bullet$ is \emph{weakly equivariant} at $x \in X$ if for all $g \in G$, $\Bar{\mu_{gx}} = g_* \Bar{\mu_x}$.
    $\mu_\bullet$ is \emph{weakly equivariant} if it is weakly equivariant at all $x \in X$.
\end{definition}

\begin{remark}
    If a weighted frame $\mu_\bullet$ is equivariant at $x\in X$, then it is weakly equivariant at $x$.
    \begin{proof}
        Let  $g \in G$ and let $E \sub G$ be a Borel set.
        Then
        \begin{align*}
            \Bar{\mu_{gx}}(E)
            &= \angles{\mu_{gx}}_{gx} (E) \\
            &= \int_{h \in G_{gx}} h_* \mu_{gx} (E) dh \\
            &= \int_{h \in G_{gx}} \mu_{gx} (\inv{h} E) dh \\
            &= \int_{h \in G_{gx}} g_* \mu_x (\inv{h} E) dh \\
            &= \int_{h \in G_{gx}} \mu_x (\inv{g} \inv{h} E) dh \\
            &= \int_{h \in gG_x \inv{g}} \mu_x (\inv{g} \inv{h} E) dh \\
            &= \int_{h \in G_x} \mu_x (\inv{h}\inv{g} E) dh \\
            &= g_* \int_{h \in G_x} h_* \mu_x(E) dh \\
            &= g_*\angles{\mu_x}_x (E) \\
            &= g_* \Bar{\mu_x} (E).
        \end{align*}
    \end{proof}
\end{remark}

\begin{remark}
    If a weighted frame $\mu_\bullet$ is weakly equivariant at $x\in X$ and $G_x$ is trivial, then it is equivariant at $x$.
    \begin{proof}
        Let  $g \in G$ and let $E \sub G$ be a Borel set.
        Then
        \begin{align*}
            \mu_{gx} (E)
            &= e_* \mu_{gx} (E) \\
            &= \int_{h \in G_x} h_* \mu_{gx} (E) dh \\
            &= \int_{h \in G_{gx}} h_* \mu_{gx} (E) dh \\
            &= \Bar{\mu_{gx}} \\
            &= g_* \Bar{\mu_x} \\
            &= g_* \int_{h \in G_x} h_* \mu_x (E) dh \\
            &= g_* \mu_x (E).
        \end{align*}
    \end{proof}
\end{remark}

\begin{proposition}
    (Proposition 3.3 in \cite{dym2024equivariant})
    Let $\mu_\bullet$ be a weakly equivariant weight frame.
    Then for any Borel function $f : X \to Y$, the function $\II_\mu (f)$ given by
    \[
        \II_\mu (f) (x) := \int_G f(\inv{g} x) d\mu_x(g)
    \]
    is $G$-invariant.
    \begin{proof}
        First, we observe
        \begin{align*}
            \II_\mu (f) (x)
            &= \int_G f(\inv{g} x) d\mu_x (g) \\
            &= \int_{G_x} \int_G f(\inv{g} x) d\mu_x (g) dh \\
            &= \int_{G_x} \int_G f(\inv{g} \inv{h} x) d\mu_x (g) dh \\
            &= \int_{G_x} \int_G f(\inv{g} x) dh_* \mu_x (g) dh \\
            &= \int_G f(\inv{g} x) d\bigp{\int_{G_x} h_* \mu_x (g) dh} \\
            &= \int_G f(\inv{g} x) d\Bar{\mu_x} (g).
        \end{align*}
        Then we simply compute
        \begin{align*}
            \II_\mu (f) (hx)
            &= \int_G f(\inv{g} h x) d\Bar{\mu_{hx}}(g) \\
            &= \int_G f(\inv{g} h x) dh_* \Bar{\mu_x}(g) \\
            &= \int_G f(\inv{g} h x) d \Bar{\mu_x}(\inv{h}g) \\
            &= \int_G f(\inv{hg} h x) d \Bar{\mu_x}(g) \\
            &= \int_G f(\inv{g} x) d \Bar{\mu_x}(g) \\
            &= \II_\mu (f) (x).
        \end{align*}
    \end{proof}
\end{proposition}

\begin{definition}
    A weakly equivariant weighted frame $\mu_\bullet$ continuous at $x$ if given any sequence $\seq{x} \sub X$ converging to $x$, we have $\angles{\mu_{x_n}}_x \wto \angles{\mu_x}_x$, i.e. for all $f : G \to \R$
    \[
        \int_{G} f d\angles{\mu_{x_n}}_x
        \to \int_{G} f d\angles{\mu_{x}}_x
        = \int_{G} f d\Bar{\mu_{x}}
    \]
    or equivalently
    \[
        \int_{G_x} \int_{G} f dh_*\mu_{x_n} dh
        \to \int_{G_x} \int_{G} f dh_*\mu_{x} dh.
    \]
    
    A \emph{robust frame} is a continuous weakly equivariant weight frame.
\end{definition}

\begin{theorem}
    (Proposition 3.6 in \cite{dym2024equivariant})
    If $\mu_\bullet$ is a robust frame for a compact group $G$ acting on a space $X$, then $\II_\mu$ preserves continuity.
    \begin{proof}
        Let $x \in X$ with $\seq{x} \sub X$ such that $x_n \to x$; we wish to show that $\II_\mu (f) (x_n) \to \II_\mu (f) (x)$ for any function $f : X \to R$.
        \begin{align*}
            \abs{\II_\mu (f) (x_n) - \II_\mu (f) (x)}
            &= \abs{\int_G f(\inv{g} x_n) d\mu_{x_n} (g) - \int_G f(\inv{g} x) d\mu_{x} (g)} \\
            &= \abs{\int_G f(\inv{g} x_n) d\mu_{x_n} (g) - \int_G f(\inv{g} x) d\angles{\mu_{x}}_x (g)} \\
            &\leq \abs{\int_G f(\inv{g} x_n) d\mu_{x_n} (g) - \int_G f(\inv{g} x) d\angles{\mu_{x_n}}_x (g)} \\
            &\quad + \abs{\int_G f(\inv{g} x) d\angles{\mu_{x_n}}_x (g) - \int_G f(\inv{g} x) d\angles{\mu_{x}}_x (g)}.
        \end{align*}
        Now observe that the latter term in this last expression has
        \[
            \abs{\int_G f(\inv{g} x) d\angles{\mu_{x_n}}_x (g) - \int_G f(\inv{g} x) d\angles{\mu_{x}}_x (g)} \to 0
        \]
        since $\mu_\bullet$ is robust, so it remains to study the former term.
        Then we have
        \begin{align*}
            \int_G f(\inv{g} x) d\angles{\mu_{x_n}}_x (g)
            &= \int_{G_x} \int_G f(\inv{g} x) d h_* \mu_{x_n} (g) dh \\
            &= \int_{G_x} \int_G f(\inv{g} \inv{h} x) d \mu_{x_n} (g) dh \\
            &= \int_{G_x} \int_G f(\inv{g} x) d\mu_{x_n} (g) dh \\
            &= \int_G f(\inv{g} x) d\mu_{x_n} (g)
        \end{align*}
        so the former term is
        \begin{align*}
            \abs{\int_G f(\inv{g} x_n) d\mu_{x_n} (g) - \int_G f(\inv{g} x) d\angles{\mu_{x_n}}_x (g)}
            &= \abs{\int_G f(\inv{g} x_n) d\mu_{x_n} (g) - \int_G f(\inv{g} x) d\mu_{x_n} (g)} \\
            &= \abs{\int_G (f(\inv{g} x_n)  - f(\inv{g} x)) d\mu_{x_n} (g)} \\
            &\leq \int_G \abs{f(\inv{g} x_n)  - f(\inv{g} x)} d\mu_{x_n} (g) \\
            &\leq \mu_{x_n}(G) \max_{g \in G} \abs{f(\inv{g} x_n) - f(\inv{g} x)} \\
            &= \max_{g \in G} \abs{f(\inv{g} x_n) - f(\inv{g} x)} \\
            &\to 0.
        \end{align*}
        Therefore we conclude that $\abs{\II_\mu (f) (x_n) - \II_\mu (f) (x)} \to 0$ and $\II_\mu (f)$ is continuous.
    \end{proof}
\end{theorem}

% TODO: compute cardinality bounds for frames in more groups/spaces

% TODO: construct frames for specific groups/spaces